% 已验证能力 DAG — 预研方法论讨论要点（草案）
% 编译：bash ../../scripts/xelatex-clean.sh 2026-07-15-已验证能力DAG预研方法论要点.tex
\documentclass[UTF8,a4paper,11pt]{ctexart}
\usepackage{amsmath,amssymb}
\usepackage{booktabs}
\usepackage{geometry}
\usepackage{hyperref}
\usepackage{longtable}
\usepackage{array}
\usepackage{listings}
\geometry{margin=2.2cm}
\newcolumntype{L}[1]{>{\raggedright\arraybackslash}p{#1}}
\hypersetup{colorlinks=true,linkcolor=blue,urlcolor=blue}
\lstset{basicstyle=\ttfamily\small,breaklines=true,frame=single}

\title{已验证能力 DAG 与预研方法论\\{\large 讨论要点草案（v0）}}
\author{ascendc 工程}
\date{2026-07-15}

\begin{document}
\maketitle

\begin{abstract}
本文档汇总近期关于「用数学模型约束预研写码路径、降低 Agent 幻觉」的讨论要点，
定位为 \texttt{docs/research/} 调研草稿，\textbf{非}定稿技术总结，\textbf{非}交付验收依据。
核心对象是域无关的\textbf{已验证能力图}及其上的\textbf{动态派生系统}；
以 ML-KEM-1024 \texttt{kem.keygen} 为样例校准，计划用 \texttt{kem.encaps} /
\texttt{kem.decaps} 等任务做规则验证实验。方法论允许随专题讨论迭代，
不指望一次提炼出最终形态。后续若干日将围绕本方向展开专题讨论，中间可穿插测试与写码。
\end{abstract}

\tableofcontents
\newpage

\section{动机与问题陈述}

\subsection{我们实际在做什么}
在 KEM KeyGen、PKE KeyGen / 加解密等路径上，开发方式是\textbf{层层递进}的：
先做最底层单点技术实验，再在已通过的能力上拓展计算任务。
一部分能力被视为「当前不可动摇的底层事实」——例如 NTT、内积、哈希——
而其下还有更底层的平台事实，例如 DataCopy、向量加、矩阵乘与 tiling、同步壳等。
后续进展默认建立在这些事实已经正确验证过的前提上。

\subsection{顶层任务的结构直觉}
顶端的 KEM 类算子，既像树状结构的根（从顶向下拆依赖），
又像 DAG 的汇点（多条已验证中间能力汇入同一目标）。
共享中间节点（如 NTT、PKE Encrypt）使结构更接近\textbf{有向无环图}而非树。

\subsection{要解决的方法问题}
能否构造一套数学模型与门禁规则，使 Agent 在开发复杂计算任务
（例如未来的 ML-DSA，或其他非密码预研）时：
\begin{itemize}
  \item \textbf{只从}已验证底层事实，以及由它们搭建并获证的中间节点出发；
  \item 高效完成更复杂任务的实现规划与落地；
  \item 从机制上压缩「跳层发明、未证书边、看似自洽却无验收」的幻觉空间。
\end{itemize}
目标是形成可指导\textbf{一般预研开发}的方法论，而不是仅服务 ML-KEM / ML-DSA 实现。

\section{方法立场（迭代原则）}

\begin{table}[h]
  \centering
  \begin{tabular}{@{}L{3.2cm}L{11cm}@{}}
    \toprule
    原则 & 含义 \\
    \midrule
    描述先于实现 & 任务启动时先产出「依赖闭包 + 缺项」，再写码 \\
    图服务门禁，不替代标准 & DAG / 自动机约束「从哪砌砖」；FIPS / liboqs 等仍是语义 oracle \\
    每验证一条边可改规则 & Encaps / Decaps 等实验若暴露漏洞，优先改方法论 \\
    粒度：可独立验收块 & 不必细到每个 API；过细则不可维护 \\
    逐渐清晰 & 允许草案演进；不一次钉死 Rule / Skill \\
    \bottomrule
  \end{tabular}
\end{table}

与仓库现有实践的关系：\texttt{baseline-registry}、活跃 \texttt{INDEX}、
\texttt{frozen} 判决、\texttt{customspec} 门禁，已经在强制
「golden / 实现路径只能引用已登记或已判决可继任的节点」。
本稿试图把隐性 DAG \textbf{显式化}，并补上「动态证书」与「合法派生」一层。

\section{对象模型（v0）}

\subsection{节点（能力事实）}
每个节点建议至少携带：

\begin{table}[h]
  \centering
  \begin{tabular}{@{}ll@{}}
    \toprule
    字段 & 说明 \\
    \midrule
    \texttt{id} & 稳定名，如 \texttt{pke.keygen.k4}、\texttt{prim.ntt.tag5t.polybatch} \\
    \texttt{layer} & \texttt{P0} 平台公理 / \texttt{P1} 密码（或领域）原语 / \texttt{P2} 构件 / \texttt{P3} 顶层算子 \\
    \texttt{iface} & I/O、dtype、形状/布局、同步接缝；\textbf{可复用的是 iface，不是源码拷贝} \\
    \texttt{cert} & 证书：路径 + CPU+SIM（及交付级 KAT）+ commit / 日期 \\
    \texttt{forbid} & 否定约束（如某阶段禁 Gather；禁 Host 密码；禁从 frozen 抄实现） \\
    \texttt{repo} & 权威锚点：\texttt{ascendc-tests/\ldots} 或 \texttt{examples/stable/\ldots} \\
    \bottomrule
  \end{tabular}
\end{table}

\subsection{边的三种类型（必须分清）}
\begin{table}[h]
  \centering
  \begin{tabular}{@{}L{2.8cm}L{6.5cm}L{5.2cm}@{}}
    \toprule
    边类型 & 含义 & KeyGen 例 \\
    \midrule
    语义依赖 & 正确性依赖下游 iface & NTT $\to$ PKE.KeyGen \\
    接缝（compose） & $A$、$B$ 各自正确，$A\!\circ\!B$ 仍须单独成节点获证 & prep $\Vert$ compute $\to$ device-keygen \\
    否定边 & 显式禁止的依赖或模板来源 & frozen 判决 $\not\to$ 活跃实现 \\
    \bottomrule
  \end{tabular}
\end{table}

\textbf{接缝不免费}：组合正确性不能由子节点证书自动推出。
这与仓库中「探针 $\to$ 集成 $\to$ stable」的实际翻车位置一致。

\subsection{证书状态}
建议状态机（可再精简）：
\[
\textsf{unproven}
\to \textsf{probe}
\to \textsf{lemma}
\to \textsf{deliverable}
;\quad
\textsf{dirty}\ \text{（上游 iface / tiling 变更后下游失效）}.
\]
\underline{lemma}：可被多父节点复用；\underline{deliverable}：stable + registry 级；
\underline{dirty}：使引理库\textbf{非单调}（可收缩），需进入动态机，而非静态证明堆。

\subsection{分层示意（以 KEM.KeyGen 为样例，非源码树）}
\begin{lstlisting}
P0  DataCopy / 向量算子 / MMAD+tiling / PIPE 同步壳
P1  NTT(poly-batch) · basemul/dot · Hash/SHAKE · CBD · ByteEncode ...
P2  Alg.13 PKE.KeyGen（及内部接缝：Â prep、se sample、compute ...）
P3  Alg.19 KEM.KeyGen = G(d,z) + Alg.13 + dk 展开 + 生产 I/O 契约
\end{lstlisting}

\section{合法性规则（Agent 门禁 v0）}

接到目标 $T$（如 \texttt{kem.encaps}）时，只允许：
\begin{enumerate}
  \item \textbf{展开} $T$ 的依赖闭包（语义边 + 接缝边）。
  \item \textbf{分类}：已有 \texttt{lemma}/\texttt{deliverable}；缺项；仅有否定边或 frozen。
  \item \textbf{缺项} $\Rightarrow$ 新建探针或 exp，\textbf{经验证写入图}后再被 $T$ 引用；
        禁止在 $T$ 任务中「顺便发明」未登记的 P0/P1。
  \item \textbf{组合}只通过已声明接缝；新接缝 $=$ 新节点，须认证。
  \item \textbf{验收}：$T$ 的 I/O 与 golden / KAT 一致；禁止以「与参考源码同构」为证。
  \item \textbf{自包含}：实现挂 iface 契约，不跨目录偷源码（\texttt{library/shared} 等仓库约定除外）；
        图边 $\neq$ \texttt{\#include} 路径。
\end{enumerate}

一句话：\textbf{只能沿已认证边生长；幻觉 $\approx$ 画出一条没有证书的边。}

\section{任务工作流（四步）}

\begin{lstlisting}
1. 目标声明  T = Alg.xx，I/O 与验收级别（probe / lemma / deliverable）
2. 闭包展开  列出 P0--P3；标 ok / missing / dirty；写出接缝清单
3. 最小补洞  只补 missing 与新接缝（先探针再集成）；否定边写入 forbid
4. 闭环写回  新节点入图；失败则改规则或封路线，不默认可复用
\end{lstlisting}

\texttt{customspec} / \texttt{STATUS} / \texttt{baseline-registry} 可视作图的\textbf{文书形态}；
方法论要求：Agent \textbf{先交闭包表，再写码}。

\section{与形式语言、自动机的关联}

\subsection{为何不止于「画一张 DAG」}
静态 DAG 适合描述\textbf{已完成引理库}的依赖投影。
预研过程是动态的：证书增减、接缝认证、frozen 抑制、dirty 失效传播。
因此需要同时借用：
\begin{itemize}
  \item \textbf{形式语言}：什么算合法的「实现计划 / 任务展开」；
  \item \textbf{自动机 / 工作流}：什么配置转移合法，何谓接受与拒绝。
\end{itemize}

\subsection{语言侧（草图）}
取字母表 $\Sigma =$「已认证叶子 iface 的调用」+ 有限组合子
（分核、分 launch、workspace 布局、host 调度等）。则：
\begin{table}[h]
  \centering
  \begin{tabular}{@{}L{3.5cm}L{5.5cm}L{5cm}@{}}
    \toprule
    语言概念 & 预研侧 & 作用 \\
    \midrule
    字 $w\in\Sigma^*$ & 一条实现计划 / 展开轨迹 & 可检查、可归档 \\
    语言 $L$ & 全体合法计划 & 「不幻觉」$\approx$ 不产出 $L$ 外的字 \\
    产生式 & 目标改写：$T\Rightarrow u_1\cdots u_k$ & 约束如何拆任务 \\
    成员判定 & 对闭包表做合法派生检查 & 动手前拦截 \\
    禁字 / 否定信息 & frozen、forbid、未登记块 & 显式「不可说」 \\
    \bottomrule
  \end{tabular}
\end{table}

$L$ 未必是正则语言：嵌套目标、接缝、属性（形状 / dtype）更接近
上下文无关或\textbf{属性文法}，或判断式 $\Gamma \vdash T$。

\subsection{自动机侧（草图）}
配置
\[
q = (\Gamma,\ F,\ S),
\]
其中 $\Gamma$ 为已证书节点集，$F$ 为前沿目标，$S$ 为接缝与约束（含 forbid）。
候选转移：
\begin{table}[h]
  \centering
  \begin{tabular}{@{}ll@{}}
    \toprule
    转移 & 含义 \\
    \midrule
    $\mathsf{Expand}$ & 将 $T\in F$ 按产生式拆成子目标 \\
    $\mathsf{Probe}$ & 为缺项建探针任务（扩展工作，不立刻进 $\Gamma$） \\
    $\mathsf{Certify}$ & 验收通过 $\Rightarrow$ 节点入 $\Gamma$ \\
    $\mathsf{Compose}$ & 仅当子节点均在 $\Gamma$ 时，接缝节点可认证 \\
    $\mathsf{Dirty}$ & 上游 iface 变 $\Rightarrow$ 下游退出 $\Gamma$ \\
    $\mathsf{Freeze}$ & 加入否定：禁止再走某边 \\
    \bottomrule
  \end{tabular}
\end{table}

\begin{itemize}
  \item \textbf{接受}：目标 $T^*\in\Gamma$ 且生产 I/O 契约满足；
  \item \textbf{拒绝阱}：使用未证书边、frozen 边、跳层发明终端。
\end{itemize}

可工作叙事（待专题讨论 refinement）：
\begin{quote}
预研方法论 $=$ 定义配置、合法转移与接受条件；\\
静态 DAG 是引理库 $\Gamma$ 的依赖投影；\\
Agent 只在该自动机上搜索接受路径；幻觉定义为非法转移。
\end{quote}

\subsection{相对古典「只增证明堆」的差异点}
\texttt{dirty} 使 $\Gamma$ \textbf{可收缩}。这一非单调性应视为方法论特征，而非瑕疵：
证书绑定具体 iface 与实现锚点，上游变更必须传播失效。

\section{域无关性：不只服务 ML-KEM}

抽出密码语义后，四类对象仍成立：

\begin{table}[h]
  \centering
  \begin{tabular}{@{}ll@{}}
    \toprule
    对象 & 一般预研含义 \\
    \midrule
    终端（axiom） & 本任务不再证明的底层事实（探针、平台行为、外部 oracle） \\
    非终端（goal） & 待构造能力 \\
    产生式（compose） & 如何用已有能力拼出新能力；接缝是一等公民 \\
    证书（judgment） & $\Gamma \vdash u:\mathsf{cert}$ \\
    \bottomrule
  \end{tabular}
\end{table}

同一套问题在非密码预研同样出现：终端从哪来、组合是否免费、失败如何进禁止规则、
Agent 是整算子生成还是合法转移搜索、进度如何用 $|\Gamma|$ / dirty 波及度量。
ML-KEM 因其 oracle 清晰、接缝多、frozen 否定边丰富，适合作\textbf{压力测试床}，
而不是方法论边界。

\section{样例校准：\texttt{kem.keygen}}

用已交付的 ML-KEM-1024 KEM KeyGen（stable + device 探针族）校准规则是否够用：

\begin{table}[h]
  \centering
  \begin{tabular}{@{}L{3.2cm}L{10.5cm}@{}}
    \toprule
    检查项 & 问题 \\
    \midrule
    分层是否干净 & P3 是否只依赖已证书 P2，而非直接重写 NTT 等 P1 \\
    接缝是否显式 & 2-launch、Alg.16 尾、\texttt{dk} 展开等是否各自有证或明确归属 \\
    registry 是否对齐 & golden 计算块 $\leftrightarrow$ 图上 lemma，有无旁路未登记内核 \\
    否定边是否写入 & Host 禁密码默认路径、禁移植参考实现为 AscendC 规格 \\
    \bottomrule
  \end{tabular}
\end{table}

缺口记为\textbf{方法论债}，进入专题讨论修改 schema；不假装图已完备。

仓库锚点（讨论时查阅，不在此复述实现）：
\begin{itemize}
  \item \texttt{examples/stable/stable-fips203-mlkem-kem-keygen-k4/}
  \item \texttt{docs/specs/fips203-mlkem1024-kem-keygen-baseline-registry.md}
  \item \texttt{ascendc-tests/pass-fix-f203-alg19-kem-keygen-device-k4/} 等活跃探针
\end{itemize}

\section{验证实验计划}

\begin{table}[h]
  \centering
  \begin{tabular}{@{}L{2.4cm}L{5.2cm}L{6.4cm}@{}}
    \toprule
    实验 & 目的 & 期望学到什么 \\
    \midrule
    A.\ Encaps & 用 v0 规则做前瞻闭包 & 是否强制挂已证 PKE.Encrypt / Hash，而非 correctness vendor；漏算接缝否 \\
    B.\ Decaps & 多父 DAG & 同时依赖 Encrypt+Decrypt+KeyGen I/O；失败路径是否成独立接缝 \\
    C.\ 对照失败 & 故意跳层描写 & 规则能否在动手前拦截 \\
    D.\ 第二域 & 非 KEM 预研链套同一机 & 证明不是密码专用话术 \\
    \bottomrule
  \end{tabular}
\end{table}

成功标准\textbf{不是}「任务一次写对」，而是：偏差能被归因到
「缺边 / 假边 / 脏证书」之一，并\textbf{回写规则}。

近期工程穿插：WSL / 探针侧 Encaps device 等写码与测试可作为实验 A 的素材，
但本专题讨论焦点仍是方法论规则，而非单次用例 PASS。

\section{建议的深挖日程（M0--M5）}

\begin{enumerate}
  \item[\textbf{M0}] 词汇表：axiom / lemma / seam / forbid / dirty / derive（域无关）。
  \item[\textbf{M1}] 用 \texttt{kem.keygen} \textbf{事后重放}一段合法派生轨迹（纸面即可）。
  \item[\textbf{M2}] 写出 Encaps / Decaps 的 $q_0$ 与允许转移；标缺项。
  \item[\textbf{M3}] 成员检查：对闭包表做合法/非法判定；故意非法计划应被拒。
  \item[\textbf{M4}] 第二域试套（如纯 tiling / 同步探针晋级链）。
  \item[\textbf{M5}] 结论稳定后，按写作模板迁入 \texttt{docs/notes/}；是否挂钩 Rule/Skill 另议。
\end{enumerate}

刻意顺序：KeyGen 重放（校准概念）$\to$ Encaps 前瞻（校准动态）$\to$ 换域（校准一般性）$\to$ 再谈工具化。

\section{刻意未决议事项}

\begin{itemize}
  \item 节点是否 $1$:$1$ 对应仓库目录名；
  \item 图用 Markdown 表还是机器可读 YAML / JSON；
  \item 何时写入 Rule / Skill（建议至少 Encaps+Decaps 各一轮后再谈）；
  \item 形式化深度：工作流网 / Petri 网 / 属性文法 / 判断式，选哪一种作为主叙事；
  \item 完整可判定性与复杂度（除非单开理论主线，否则暂不作为预研交付门槛）。
\end{itemize}

\section{立论草案（一句话）}

\begin{quote}
预研开发的合法性，可刻画为：在「已证书能力」生成的语言 $L$ 上，
由配置自动机执行 \textsf{Expand}/\textsf{Probe}/\textsf{Certify}/\textsf{Compose}/\textsf{Dirty}/\textsf{Freeze}；
静态 DAG 是引理库 $\Gamma$ 的依赖投影，不是全过程。
Agent 的职责是在该自动机上搜索接受路径；幻觉被定义为非法转移。
\end{quote}

\section{文档地位与后续}

\begin{itemize}
  \item 路径：\texttt{docs/research/} —— 调研草稿，可废弃、可改写。
  \item 本文件记录的是\textbf{方法论讨论要点}，不记录题外话的保护策略讨论。
  \item 专题讨论期间允许穿插测试与写码；结论晋级 \texttt{notes/} 前，不以本 PDF 约束交付验收。
  \item 维护：增删条目同步 \texttt{docs/research/INDEX.md}；重大决策同步当日 \texttt{qa/}。
\end{itemize}

\end{document}
